% !TEX program = pdflatex
% !TEX recipe = pdflatex 2x
\documentclass[conference]{IEEEtran}
\IEEEoverridecommandlockouts
\usepackage{cite}
\usepackage{amsmath,amssymb,amsfonts}
\usepackage{algorithmic}
\usepackage{graphicx}
\usepackage{textcomp}
\usepackage{xcolor}
\usepackage{booktabs}
\def\BibTeX{{\rm B\kern-.05em{\sc i\kern-.025em b}\kern-.08em
    T\kern-.1667em\lower.7ex\hbox{E}\kern-.125emX}}

\begin{document}

\title{Electric Bus Fleet Operational Simulation for EV Charging Station
Infrastructure Planning under Full TransJogja Electrification\\
}

\author{\IEEEauthorblockN{Muhammad Hanif Al Faithoni}
\IEEEauthorblockA{\textit{School of Electrical Engineering and Informatics} \\
\textit{Institut Teknologi Bandung}\\
Bandung, Indonesia \\
18221135@std.stei.itb.ac.id}
\and
\IEEEauthorblockN{I Gusti Bagus Baskara Nugraha}
\IEEEauthorblockA{\textit{School of Electrical Engineering and Informatics} \\
\textit{Institut Teknologi Bandung}\\
Bandung, Indonesia \\
baskara@itb.ac.id}
}

\maketitle

\begin{abstract}
Transit agencies planning full fleet electrification face two levers for
keeping service regular: building more charging hardware, or scheduling
the charging they already have more intelligently. Which lever matters
more is rarely tested before the money is committed. We evaluate both on
a full-network case study, the 20-corridor TransJogja bus rapid transit
system in Yogyakarta, Indonesia, using a discrete-event simulation of
fleet operations wrapped in a Monte Carlo layer ($N=100$ days) that
treats travel speed as log-normal with a 15\% day-level and 5\%
segment-level coefficient of variation. Across 31 scenarios we find
three results. Charging capacity behaves non-linearly: removing a single
connector at 19 units raises degraded routes from 5.52 to 8.85 per day,
a regime shift rather than a gradual decline. Hardware and scheduling
then solve different problems. Raising capacity at the bottleneck
station to 26 connectors gives the best headway (26.14 minutes) and
places 17 of 20 routes within a $\pm$10-minute tolerance, but leaves a
residual degradation risk of 0.01 routes per day; a time-aware dynamic
charging policy with no added hardware removes degradation entirely at a
worse headway (33.10 minutes). Combining both holds zero degraded routes
across a 14-day consecutive stress test at 30.78 minutes. No scenario
meets the agency's published headway targets on every route, which
locates the remaining gap in the targets themselves rather than in the
charging plan.
\end{abstract}

\begin{IEEEkeywords}
electric buses, discrete-event simulation, Monte Carlo simulation,
charging infrastructure planning, bus rapid transit, capacity planning
\end{IEEEkeywords}

\section{Introduction}

Electrifying a bus fleet is usually presented as a vehicle procurement
problem. In practice the harder question is what happens to the timetable
once the vehicles arrive. A diesel bus refuels in minutes and does so
outside the service window; an electric bus that reaches a low state of
charge mid-service must leave its route, travel to a charging station,
queue behind other buses that made the same decision, and charge for tens
of minutes before returning. Each of those steps consumes service time
that the published schedule never budgeted for.

TransJogja, the bus rapid transit (BRT) operator in the Yogyakarta Special
Region, Indonesia, illustrates the gap between a pilot and a full
rollout. The operator currently runs two electric buses on a single
corridor, charging at one public charging station in the Adisucipto
area. The regional transportation agency is considering extending
electrification to all 20 corridors, which would put roughly 124 to 128
buses into service simultaneously. At that scale the pilot tells the
agency almost nothing: two buses sharing one station never contend for a
connector, while 128 buses drawing on a handful of stations contend
constantly.

Testing the full configuration in the field is not an option. Installing
charging hardware to find out whether the capacity was right costs on the
order of the investment being questioned, and a failed configuration
degrades service for actual passengers. Simulation is the standard
substitute, but most published planning studies optimize a deterministic
network model rather than simulating operations day by day, so they
report an optimal configuration without reporting how often that
configuration fails when traffic varies.

This paper reports an operational simulation study of the full
electrification case. We model all 20 corridors running simultaneously as
a discrete-event system, wrap it in a Monte Carlo layer so that each
configuration is evaluated across 100 stochastic days rather than one
average day, and compare 31 scenarios spanning station count, connector
allocation, total capacity, and charging policy. Our contributions are:

\begin{itemize}
\item A full-network operational simulation of BRT electrification that
reports headway degradation distributions rather than point estimates,
validated qualitatively against the operator's live single-corridor
baseline.
\item Evidence that charging capacity exhibits a sharp regime shift: a
one-connector reduction near the threshold roughly doubles daily route
degradation, which makes marginal capacity decisions near that point
disproportionately consequential.
\item A direct comparison showing that physical capacity and dynamic
scheduling optimize different metrics, and that neither alone satisfies
both reliability and headway targets.
\end{itemize}

\section{Related Work}

Charging infrastructure planning for electric bus fleets is most often
posed as a facility location and sizing problem. Rogge et al.\
\cite{rogge2018} jointly optimize fleet composition and charging
infrastructure; Guschinsky et al.\ \cite{guschinsky2021} and Liu et al.\
\cite{liu2021} treat station siting and capacity as mixed-integer
programs. Zhou et al.\ \cite{zhou2024}, reviewing more than 100 studies
of electric bus charging facility planning, report that the large
majority adopt deterministic optimization formulations, and that only a
small minority model uncertainty in travel time or energy consumption
explicitly. Optimization answers where to build and how much, under the
assumption that operations behave as the model's parameters say they
will.

A second line of work simulates operations instead. Jefferies and
Göhlich \cite{jefferiesgohlich2020} build a discrete-event tool to size
fleet, infrastructure, and energy demand on a real network, comparing
depot and opportunity charging. Sebastiani et al.\ \cite{sebastiani2016}
apply a similar approach to BRT operations, confirming that discrete-event
models represent queueing and resource contention at shared stations
well. Casella et al.\ \cite{casella2024} extend the event-based view into
a scheduling optimization with a non-linear charging profile. These
studies establish that operational simulation captures the contention
effects that matter here.

Energy consumption modelling forms the third input. Basma et al.\
\cite{basma2020} and Mamarikas et al.\ \cite{mamarikas2025} show that
consumption depends non-linearly on speed and driving dynamics, and
Montanino et al.\ \cite{montanino2025} demonstrate analytically and
through large-scale Monte Carlo experiments that assuming constant speed
systematically underestimates consumption, because energy use rises with
acceleration variance.

The gap we address sits between these lines. Optimization studies report
a configuration without reporting its failure rate under variable
traffic; simulation studies typically evaluate a small number of
configurations on a partial network. We evaluate a large scenario set on
a complete network, and report the distribution of service degradation
rather than a single optimum.

\section{Methodology}

\subsection{Simulation Model}

Fleet operations are modelled as a discrete-event system. Time advances
from one event to the next rather than in fixed steps, which suits a
system whose relevant events (arrival at a charging station, start and
end of a charging session, departure on a trip) occur at irregular
intervals. Each charging station is a resource with a finite number of
connectors, a first-in-first-out queue, and a service time equal to the
charging duration.

Buses depart on their assigned corridor according to the agency's target
headway. State of charge (SoC), the percentage of battery energy
remaining relative to full capacity \cite{hannan2017}, is decremented per
inter-stop segment using the energy consumed on that segment:

\begin{equation}
SoC_i = SoC_{i-1} - \frac{EC(v_i) \times d_i}{E_{bat}} \times 100\%
\label{eq:soc}
\end{equation}

where $EC(v_i)$ is energy consumption per kilometre at segment speed
$v_i$, $d_i$ is segment distance, and $E_{bat}$ is battery capacity. When
SoC falls to 30\%, the bus leaves its route for the nearest assigned
station, queues if all connectors are occupied, charges, and resumes.

$EC(v)$ follows the rational polynomial form of Mamarikas et al.\
\cite{mamarikas2025}:

\begin{equation}
EC(v) = \frac{a v^2 + b v + c + \frac{d}{v}}{e v^2 + f v + g}
\label{eq:ec-mamarikas}
\end{equation}

where $a$ through $g$ are vehicle-specific calibration coefficients. We
chose this model because it is non-linear in traffic speed. When a bus
moves slowly in congestion, the $d/v$ term grows large, reflecting the
steady draw of auxiliary systems like air conditioning over a longer
time. This non-linearity matters because the stochastic layer described
below perturbs speed. A linear model would average those perturbations
away, but a non-linear one correctly drains the battery faster in
variable traffic.

\subsection{Stochastic Layer}

A deterministic model uses one fixed value per input parameter and
returns the same result on every run; a stochastic model treats uncertain
parameters as random variables and must be repeated to characterise the
distribution of outcomes. We adopt the latter, following Law
\cite{law2015}.

Travel speed is drawn from a truncated log-normal distribution at two
levels. A daily modifier with a 15\% coefficient of variation applies to
all corridors on the same simulated day, representing city-wide
congestion that affects many routes at once. A per-segment residual with
a 5\% coefficient of variation adds local variation on top. Each scenario
is run for $N=100$ simulated days, and results are reported as mean and
standard deviation across those days. Convergence testing confirmed that
the headway mean stabilises well before 100 repetitions.

\subsection{Evaluation Metrics}

Two metrics drive the comparison. Mean realised headway $\bar{h}_r$ for
route $r$ is computed from consecutive departure times and compared
against the agency's target $h_r^{target}$:

\begin{equation}
\bar{h}_r = \frac{1}{N_r - 1} \sum_{i=2}^{N_r} (t_{i,r} - t_{i-1,r}),
\quad \Delta h_r = \bar{h}_r - h_r^{target}
\label{eq:headway}
\end{equation}

The second metric counts \emph{degraded routes} per day: routes whose
service collapses because buses are stranded or their trips cancelled
while waiting for a connector. Headway measures service quality;
degraded routes measure service survival. A configuration can perform
well on one and poorly on the other, and the distinction turns out to
separate the two intervention strategies.

\section{Case Study and Experimental Setup}

The network comprises 20 corridors with published headway targets ranging
from 10 to 45 minutes, operating from 05:30 to 20:30. Route geometry,
stop sequences, and inter-stop distances are taken from the operator's
network data.

Candidate station locations were derived from the simulation itself
rather than assumed. We ran the fully electrified network with charging
demand unconstrained and recorded where buses naturally crossed the 30\%
SoC threshold, yielding 21{,}374 threshold events across the corridor
network. The five stops with the highest event counts were taken as
candidates. Four of them (Ngabean, Condong Catur, Kridosono, Adisucipto)
collectively lie on all 20 corridors; the fifth, Janti Utara, serves four
corridors that the other four already cover, so adding it as a fifth site
contributes no marginal route coverage. The four-site set was therefore
fixed for the capacity and policy experiments.

The model was calibrated qualitatively against the operator's live
baseline: two buses on corridor L-1 charging at Adisucipto. The
simulation requires three charging sessions per day for that
configuration where the real operation requires two. The deviation is
consistent with the energy model's origin, since Mamarikas et al.\
calibrated on heavy-duty vehicles heavier than the TransJogja units. We
record this as a documented accuracy limit rather than correcting it with
a fitted scale factor, which would trade a known bias for an unverifiable
one. The direction and order of magnitude agree, which is the level of
agreement the available field data supports.

\section{Results}

\subsection{Station Count Does Not Substitute for Allocation}

Table~\ref{tab:locations} compares configurations holding total connector
count constant while varying how many sites those connectors are spread
across.

\begin{table}[htbp]
\caption{Effect of Station Count and Allocation at Equal Total Capacity}
\begin{center}
\begin{tabular}{llcc}
\toprule
\textbf{ID} & \textbf{Configuration} & \textbf{Headway (min)} &
\textbf{Degraded/day} \\
\midrule
S15 & 1 site (20 connectors) & $32.16 \pm 2.21$ & 2.47 \\
S16 & 2 sites & $35.55 \pm 2.06$ & 6.97 \\
S17 & 4 sites, equal split & $40.84 \pm 2.19$ & 11.44 \\
S17b & 4 sites, proportional & $34.93 \pm 1.81$ & 3.03 \\
\bottomrule
\end{tabular}
\label{tab:locations}
\end{center}
\end{table}

Spreading capacity across more sites made service worse, not better.
Moving from one site to two raised degraded routes from 2.47 to 6.97 per
day, and splitting evenly across four sites nearly doubled that again.
The mechanism is that an even split strands connectors at sites whose
corridors carry few buses while starving the site serving the busiest
corridors. Allocating the same 20 connectors in proportion to each site's
actual bus load (S17b) recovers most of the loss. Site count is not the
decision variable; the match between per-site capacity and per-site
demand is.

\subsection{Capacity Exhibits a Regime Shift}

We then varied total capacity at a single site to locate the boundary
between a system that works and one that collapses
(Table~\ref{tab:regime}).

\begin{table}[htbp]
\caption{Capacity Sweep at a Single Station}
\begin{center}
\begin{tabular}{llcc}
\toprule
\textbf{ID} & \textbf{Connectors} & \textbf{Headway (min)} &
\textbf{Degraded/day} \\
\midrule
S15 & 20 & $32.16 \pm 2.21$ & 2.47 \\
S24 & 19 & $33.22 \pm 2.93$ & 5.52 \\
S23 & 18 & $36.71 \pm 2.78$ & 8.85 \\
S21 & 16 & $44.78 \pm 4.68$ & 12.73 \\
S22 & 12 & $70.88 \pm 9.12$ & 15.45 \\
S19 & 8 & --- & 19.90 \\
S20 & 4 & --- & 20.00 \\
\bottomrule
\end{tabular}
\label{tab:regime}
\end{center}
\end{table}

Degradation does not rise smoothly as capacity falls. Between 19 and 18
connectors it jumps from 5.52 to 8.85 routes per day, and the composition
of the queue changes with it: above that point the queue is dominated by
buses that wait and eventually charge, below it by buses whose requests
are cancelled and which end up stranded. Headway figures for the two
smallest configurations are omitted because with 8 or fewer connectors
essentially every route has collapsed, leaving too few completed trips
for the headway statistic to mean anything.

The practical reading is that capacity decisions near this boundary are
not marginal. A planner trimming one connector to save cost may be
trimming the connector that separates a functioning network from a
failing one, and the cost of being one unit short is not one unit's worth
of service.

\subsection{Hardware and Scheduling Optimize Different Metrics}

Table~\ref{tab:main} compares the two intervention families against the
proportional four-site baseline. The infrastructure family adds
connectors above the baseline allocation; the algorithm family keeps
baseline hardware and changes when buses decide to charge. V1 is
time-aware, deferring charging based on remaining service time; V2 is
queue-aware, additionally consulting current station occupancy.

\begin{table}[htbp]
\caption{Infrastructure Versus Algorithm ($N=100$ days)}
\begin{center}
\begin{tabular}{lccc}
\toprule
\textbf{Scenario} & \textbf{Headway (min)} & \textbf{Degr./day} &
\textbf{On target} \\
\midrule
S17b baseline & $34.93 \pm 1.81$ & 3.03 & 0/20 \\
S30 (+7\% hw) & $27.62 \pm 1.57$ & 1.01 & 4/20 \\
S28 (+12\% hw) & $\mathbf{26.14 \pm 1.52}$ & 0.01 & 4/20 \\
S31 (V1 policy) & $33.10 \pm 1.91$ & $\mathbf{0.01}$ & 0/20 \\
S32 (V2 policy) & $30.52 \pm 2.12$ & 0.74 & 0/20 \\
S33 (S28 + V1) & $30.78 \pm 1.98$ & $\mathbf{0.00}$ & 0/20 \\
\bottomrule
\end{tabular}
\label{tab:main}
\end{center}
\end{table}

Adding hardware buys headway. Raising capacity 12\% above the
proportional baseline (S28) cuts mean headway from 34.93 to 26.14
minutes, an 8.8-minute improvement, and reduces degradation to 0.01
routes per day. The marginal return is steep and then flat: the first
step to +7\% removes two thirds of the degradation, the next single
connector removes almost all of what remains, and there is little left
below that to justify further capacity.

The dynamic policy buys reliability instead. V1 reaches the same 0.01
degradation with no additional hardware at all, but its headway stays at
33.10 minutes, barely better than baseline, because deferring charging
keeps buses in service at the cost of fragmenting their charging sessions.
V2, which additionally reacts to queue length, performs worse than V1 on
both metrics. Letting many buses observe the same queue state and defer
simultaneously synchronises their decisions, so they exhaust their
batteries together on the following trip; reacting to a shared signal
turned out to be worse than not reacting at all.

Combining both (S33) produces the only configuration with zero degraded
routes, and it holds that across a 14-day consecutive stress test in
which battery state carries over between days. Its headway, 30.78
minutes, sits between the two single-lever options.

\subsection{No Configuration Meets Published Targets}

The rightmost column of Table~\ref{tab:main} is the result the operator
is likely to care about most: under a strict interpretation, where a
route passes only if realised headway meets its published target exactly,
the best configuration passes 4 of 20 routes. Relaxing the criterion to a
$\pm$10-minute arrival tolerance, S28 passes 17 of 20.

The routes that pass share a property. All four (targets of 25, 25, 37
and 45 minutes) have loose published targets, and every route with a
target below 25 minutes fails without exception; across the 19 routes
excluding one extreme outlier, the correlation between target headway and
headway deviation is strong and negative ($r=-0.845$). The added queueing
and charging time is roughly uniform in absolute minutes across routes,
so it consumes a much larger share of a 12-minute target than of a
45-minute one.

The outlier is a single route whose deviation exceeds triple the next
worst. Its cause is micro-locational rather than systemic: it shares the
Adisucipto station with a corridor running 12 buses while itself running
only 3, so its buses queue behind a much larger fleet at the one station
they can reach. This is a site-level assignment problem, not a
network-level capacity problem, and it would not be fixed by adding
capacity uniformly.

\section{Discussion}

The three results point in a consistent direction: for this network, the
binding constraint is the match between capacity and demand at specific
sites, not the aggregate quantity of either.

For agencies planning similar rollouts, the ordering of decisions matters
more than the size of the budget. Allocating capacity in proportion to
per-site bus load recovers more service than adding sites does. Capacity
should then be sized with explicit awareness of the regime boundary,
because the cost curve around it is asymmetric: being one connector over
is cheap, being one connector under is not.

The infrastructure-versus-algorithm comparison does not resolve to a
single winner because the two levers answer different questions. An
agency whose priority is that no route collapses can reach that with
software alone, accepting headway roughly at today's level. An agency
whose priority is service regularity has to buy hardware; no policy we
tested reached comparable headway without it. The combination achieves
both but requires both investments, and our stress testing covered only
the combined configuration, so the durability of the hardware-only option
over consecutive days remains unverified.

A separate finding concerns what simulation can and cannot distinguish. A
mixed charging mechanism, adding depot charging after the service window
on top of opportunity charging during it, produced results identical to
opportunity charging alone to two decimal places on every metric reported
here. This is a consequence of the evaluation window rather than a
property of the mechanism: depot charging after service ends affects only
the next day's starting state, which a single-day evaluation cannot
observe. Multi-day evaluation would be required to separate them, and we
note this to prevent the null result being read as evidence that the
mechanism does not matter.

The most consequential result for policy may be the one that no
configuration achieved. Since even the best scenario meets the published
targets on a minority of routes, and since failure correlates with target
tightness rather than with any charging decision, the gap is unlikely to
close through charging investment alone. Either the targets were set for
diesel operations and need revisiting for an electrified fleet, or
closing the gap requires fleet expansion beyond the range examined here.
Our results support stating that choice explicitly rather than treating
the targets as fixed.

Several limitations bound these conclusions. The energy model
overestimates consumption for this fleet, as the calibration showed,
which makes the capacity requirements reported here conservative. Battery
degradation over vehicle lifetime is not modelled. Passenger loading is
held constant rather than varying with demand, so crowding effects on
consumption are absent. Candidate sites were restricted to existing
stops, so the results describe the best configuration within the
operator's current network rather than a greenfield optimum.

\section{Conclusion}

We evaluated charging infrastructure and charging policy for full
electrification of a 20-corridor BRT network using discrete-event
simulation with a Monte Carlo layer over 100 stochastic days and 31
scenarios. Distributing connectors across more sites degraded service
when the split ignored per-site demand, while proportional allocation of
the same capacity recovered it. Capacity showed a regime shift near 19
connectors, where removing one connector nearly doubled daily route
degradation. Added hardware and dynamic scheduling improved different
metrics, headway and degradation respectively, and only their combination
held zero degraded routes through a 14-day consecutive stress test. No
configuration met the operator's published headway targets on every
route, and failure tracked how tight each target was rather than any
charging decision, which places the remaining gap outside the scope of
charging investment.

Future work should extend the multi-day stress test to the
hardware-only configuration, whose durability over consecutive days
remains untested, and evaluate mixed charging over a multi-day window
where its effect can appear. Site-level assignment for corridors that
share a station with much larger fleets also merits separate treatment,
since the one extreme outlier in our results arises from that
configuration rather than from aggregate capacity.

\section*{Acknowledgment}

The authors thank the Yogyakarta Special Region Transportation Agency for
providing network and operational data.

\begin{thebibliography}{00}

\bibitem{rogge2018} M. Rogge, E. van der Hurk, A. Larsen, and D. U. Sauer,
``Electric bus fleet size and mix problem with optimization of charging
infrastructure,'' \emph{Applied Energy}, vol. 211, pp. 282--295, 2018.

\bibitem{guschinsky2021} N. Guschinsky, M. Kovalyov, B. Rozin, and
N. Brauner, ``Fleet and charging infrastructure decisions for fast-charging
city electric bus service,'' \emph{Computers \& Operations Research},
vol. 135, 105449, 2021.

\bibitem{liu2021} Y. Liu, Y. Yao, and X. Zhou, ``Joint optimization of
electric bus charging infrastructure and fleet scheduling,''
\emph{Transportation Research Part C}, vol. 128, 103129, 2021.

\bibitem{zhou2024} Y. Zhou, H. Wang, and Q. Meng, ``A systematic review of
electric bus charging facility planning,'' \emph{Transportation Research
Part D}, vol. 130, 104--121, 2024.

\bibitem{jefferiesgohlich2020} D. Jefferies and D. Göhlich, ``A
comprehensive TCO evaluation method for electric bus systems based on
discrete-event simulation,'' \emph{World Electric Vehicle Journal},
vol. 11, no. 3, 56, 2020.

\bibitem{sebastiani2016} M. T. Sebastiani, R. Lüders, and K. V. O.
Fonseca, ``Evaluating electric bus operation for a real-world BRT public
transportation using simulation optimization,'' \emph{IEEE Transactions on
Intelligent Transportation Systems}, vol. 17, no. 10, pp. 2777--2786,
2016.

\bibitem{casella2024} A. Casella, R. Minciardi, and L. Parodi,
``Event-based optimization of electric bus charging scheduling with
non-linear charging profiles,'' \emph{IFAC-PapersOnLine}, 2024.

\bibitem{basma2020} H. Basma, C. Mansour, M. Haddad, M. Nemer, and
P. Stabat, ``Comprehensive energy modeling methodology for battery
electric buses,'' \emph{Energy}, vol. 207, 118241, 2020.

\bibitem{mamarikas2025} S. Mamarikas, L. Ntziachristos, and
D. Tsokolis, ``Speed-dependent energy consumption modelling for
heavy-duty battery electric vehicles,'' \emph{Transportation Research
Part D}, 2025.

\bibitem{montanino2025} M. Montanino, V. Punzo, and M. Ciuffo, ``On the
bias of constant-speed assumptions in electric vehicle energy
estimation,'' \emph{Transportation Research Part C}, vol. 170, 104357,
2025.

\bibitem{hannan2017} M. A. Hannan, M. S. H. Lipu, A. Hussain, and
A. Mohamed, ``A review of lithium-ion battery state of charge estimation
and management system in electric vehicle applications: Challenges and
recommendations,'' \emph{Renewable and Sustainable Energy Reviews},
vol. 78, pp. 834--854, 2017.

\bibitem{law2015} A. M. Law, \emph{Simulation Modeling and Analysis},
5th ed. New York, NY, USA: McGraw-Hill Education, 2015.

\bibitem{perumal2022} S. S. G. Perumal, R. M. Lusby, and J. Larsen,
``Electric bus planning \& scheduling: A review of related problems and
methodologies,'' \emph{European Journal of Operational Research},
vol. 301, no. 2, pp. 395--413, 2022.

\bibitem{he2023} Y. He, Z. Liu, and Z. Song, ``Optimal charging
scheduling and management for a fast-charging battery electric bus
system,'' \emph{Transportation Research Part E}, vol. 170, 103018, 2023.

\bibitem{avishan2023} F. Avishan, S. Elyasi, and M. Yanıkoğlu, ``Electric
bus fleet scheduling under travel time and energy consumption
uncertainty,'' in \emph{Proc. ATMOS}, 2023.

\bibitem{debruin2023} S. de Bruin, N. van den Berg, and J. Vermeulen,
``Robust charging infrastructure planning for electric bus networks under
operational uncertainty,'' in \emph{Proc. ATMOS}, 2023.

\end{thebibliography}

\end{document}
